\documentclass[11pt]{article}

\usepackage{amsmath,amssymb}
\usepackage{xurl}
\usepackage[hidelinks]{hyperref}
\setlength{\emergencystretch}{2em}

% Shared commands for notes in docs/paper-gaps/.

\DeclareUnicodeCharacter{2080}{\ifmmode{}_{0}\else\textsubscript{0}\fi}
\DeclareUnicodeCharacter{2081}{\ifmmode{}_{1}\else\textsubscript{1}\fi}
\DeclareUnicodeCharacter{2082}{\ifmmode{}_{2}\else\textsubscript{2}\fi}
\DeclareUnicodeCharacter{2099}{\ifmmode{}_{n}\else\textsubscript{n}\fi}

\newcommand{\Lean}{\textsc{Lean}}
\ExplSyntaxOn
\NewDocumentCommand{\leanid}{m}{
  \ifmmode
    \textsf{\detokenize{#1}}
  \else
    \group_begin:
    \urlstyle{sf}
    \tl_set:Nn \l_tmpa_tl {#1}
    \tl_replace_all:Nnn \l_tmpa_tl {\_} {_}
    \exp_args:NV \path \l_tmpa_tl
    \group_end:
  \fi
}
\ExplSyntaxOff
\providecommand{\tr}{\operatorname{tr}}
\newcommand{\tnleanIssueBase}{https://github.com/LionSR/TNLean/issues/}
\newcommand{\tnleanPrBase}{https://github.com/LionSR/TNLean/pull/}
\newcommand{\tnleanDocBase}{https://sirui-lu.com/TNLean/docs/}
\newcommand{\ghissue}[1]{\href{\tnleanIssueBase#1}{GitHub issue \##1}}
\newcommand{\ghpr}[1]{\href{\tnleanPrBase#1}{PR \##1}}
\newcommand{\doclink}[1]{\href{\tnleanDocBase#1.html}{\path{#1}}}

% Dirac notation and the identity operator, as in blueprint/src/macros/common.tex.
% \providecommand keeps note-local definitions authoritative.
\usepackage{dsfont}
\providecommand{\ket}[1]{|#1\rangle}
\providecommand{\bra}[1]{\langle #1|}
\providecommand{\braket}[2]{\langle #1 | #2 \rangle}
\providecommand{\ketbra}[2]{|#1\rangle\!\langle #2|}
\providecommand{\Id}{\mathds{1}}
\providecommand{\one}{\mathds{1}}
\providecommand{\C}{\mathbb{C}}
\providecommand{\MN}[1]{M_{#1}(\C)}
\providecommand{\MD}{M_D(\C)}

% External referencing for paper sources.  The build helper writes sanitized
% auxiliary files under build/paper-aux/ and excludes duplicated source labels.
\usepackage{xr}

% Import a generated paper auxiliary file with a caller-supplied label prefix.
% The candidate paths support builds from docs/paper-gaps/, the repository root,
% and build/paper-aux/ itself.
\newcommand{\importpaperaux}[2]{%
  \IfFileExists{../../build/paper-aux/#2.aux}{%
    \externaldocument[#1]{../../build/paper-aux/#2}%
  }{%
    \IfFileExists{build/paper-aux/#2.aux}{%
      \externaldocument[#1]{build/paper-aux/#2}%
    }{%
      \IfFileExists{#2.aux}{%
        \externaldocument[#1]{#2}%
      }{}%
    }%
  }%
}

% General external reference macros
\newcommand{\paperref}[2]{\ref{#1-#2}}
\newcommand{\papereqref}[2]{(\ref{#1-#2})}

% CPSV16 source (1606.00608 / MPDO-22-12-17-2.tex) external document and shortcuts
\importpaperaux{cpsv16-}{1606.00608/MPDO-22-12-17-2-tnlean}
\newcommand{\cpsvref}[1]{\paperref{cpsv16}{#1}}
\newcommand{\cpsveqref}[1]{\papereqref{cpsv16}{#1}}

% Verdict marker: \gapnote{<kind>}{<status>}. Kind is the policy
% classification (clarification, local-correction, scope-restriction,
% unfaithful, false-source, open-gap); status is open, wip (elimination
% underway), resolved, or historical. Machine-read by the site index;
% prints a verdict line.
\newcommand{\gapnote}[2]{\par\noindent\textbf{Verdict:} #1 (#2).\par}


\title{Conditional After-Blocking Fundamental Theorem versus the CPSV Statement}
\author{}
\date{2026-05-01}

\begin{document}
\maketitle

\section{The source statement}

The non-periodic matrix product vector Fundamental Theorem in
\cite{Cirac2016MPDO_arXiv,Cirac2017MPDO_AnnPhys,Cirac2021Matrix} is not stated
with the auxiliary hypotheses used in the conditional Lean theorem. It is stated
after the tensors have already been put in canonical form and after bases of
normal tensors have been chosen. This note records that distinction for the
after-blocking theorem in \ghpr{1104}, which is part of the discussion in
\ghissue{1093}.

In the 2016 source, the canonical-form discussion first removes invariant
subspaces. After finitely many steps the matrices are replaced by block-diagonal
matrices
\[
    A^i = \bigoplus_{k=1}^r \mu_k A_k^i,
    \qquad
    \mu_k\ne0,
    \qquad |\mu_k|\le 1,
    \qquad \max_k|\mu_k|=1,
    \tag{CF}
\]
which generate the same family of matrix product vectors as the original tensor;
zero blocks are allowed in the sense that $\sum_k D_k\le D$
\cite[lines~214--219]{Cirac2016MPDO_arXiv}. The later BNT comparison arranges
the surviving basis blocks by strictly decreasing moduli, a separate ordering
convention:
\[
    |\mu_1|=1,\qquad |\mu_1|>|\mu_2|>\cdots.
    \tag{BNT-order}
\]
Periodic peripheral components are
then removed by blocking a common multiple of their periods
\cite[lines~227--231]{Cirac2016MPDO_arXiv}, and the paper records the resulting
canonical-form existence statement: after blocking, any tensor can be represented
by a tensor in canonical form generating the same family
\cite[lines~249--251]{Cirac2016MPDO_arXiv}. The review gives the same two-step
account: first canonical form with the same family of matrix product vectors, and
then the Fundamental Theorem for tensors already in canonical form
\cite[lines~1752--1758]{Cirac2021Matrix}. It also describes the direct-sum
canonical form and the period-removing blocking
\cite[lines~1798--1820]{Cirac2021Matrix}.

The comparison theorem is expressed through bases of normal tensors. In the
2016 source a basis of normal tensors is a minimal family of normal tensors such
that the matrix product vectors of the original tensor are linear combinations of
those basis vectors and the latter are eventually linearly independent
\cite[lines~271--280]{Cirac2016MPDO_arXiv}. For a tensor in canonical form, the
paper then writes the matrix product vectors as sums of the basis-of-normal-tensor
vectors with coefficients $\sum_q \mu_{j,q}^N$
\cite[lines~283--302]{Cirac2016MPDO_arXiv}. The review states the same basis
notion and the same expansion
\cite[lines~1846--1885]{Cirac2021Matrix}.
Let $a_j(N)$ and $b_\ell(N)$ denote the length-$N$ matrix product vectors
of the $j$-th and $\ell$-th basis normal tensors on the two sides. In the
notation of the BNT comparison, this is the expansion
\[
    V_A^{(N)} = \sum_{j=1}^{g_A}
    \left(\sum_{q=1}^{r_{A,j}}\mu_{A,j,q}^N\right)a_j(N),
    \qquad
    V_B^{(N)} = \sum_{\ell=1}^{g_B}
    \left(\sum_{q=1}^{r_{B,\ell}}\nu_{B,\ell,q}^N\right)b_\ell(N).
    \tag{BNT}
\]

Injectivity is supplied separately. The 2016 source defines injectivity and
block-injective canonical form, invokes the quantum Wielandt theorem to say that
normal tensors become injective after blocking, and records the block-injective
bound for canonical-form tensors
\cite[lines~317--344]{Cirac2016MPDO_arXiv}. The review likewise states that
after blocking at most $3D^5$ sites a canonical-form tensor acquires
block-injective canonical form
\cite[lines~2114--2124]{Cirac2021Matrix}.
Schematically, if $A_k$ is a normal block, the paper uses a blocking length
$L$ for which
\[
    \operatorname{span}
    \{A_k^{i_1}\cdots A_k^{i_L}:1\le i_1,\ldots,i_L\le d\}
    =
    M_{D_k}(\mathbb C).
    \tag{Inj}
\]

The CPSV statement itself is then short. For tensors $A$ and $B$ in
canonical form, with bases of normal tensors $A_{k_a}$ and $B_{k_b}$,
proportionality of all matrix product vector families implies equality of the
number of basis tensors and a matching of basis tensors up to permutation,
phases, and invertible gauge matrices
\cite[Theorem~\texttt{thm1}, lines~347--352]{Cirac2016MPDO_arXiv}. In the equal
case the paper states the corresponding global gauge conclusion
\cite[Corollary~\texttt{II\_cor2}, lines~354--360]{Cirac2016MPDO_arXiv}. The
review repeats these two statements
\cite[Theorem~\texttt{thm1} and Corollary~\texttt{II\_cor2}, lines~1891--1900]{Cirac2021Matrix}.
The review also points out that a later periodic Fundamental Theorem avoids the
period-removing blocking step
\cite[lines~1908--1909]{Cirac2021Matrix}; that theorem is intentionally not used
in the derivation considered here.
In formula form, the proportional statement is
\[
    \begin{aligned}
         & \forall N,\quad \exists c_N\in\mathbb C,\ c_N\ne0,\quad
        V_A^{(N)}=c_NV_B^{(N)}                                     \\
         & \hspace{4em}\Longrightarrow\quad
        g_A=g_B
        \quad\land\quad
        B_{\sigma(j)}^i=e^{i\phi_j}X_jA_j^iX_j^{-1}.
    \end{aligned}
    \tag{FT}
\]
The equal statement then assembles these block gauges into a global gauge, with
the finite-multiplicity version carrying identity factors on multiplicity
spaces:
\[
    Y=\bigoplus_j I_{r_j}\otimes X_j.
    \tag{Gauge}
\]

\section{The conditional Lean theorem}

The theorem added in \ghpr{1104} starts from two tensors $A$ and $B$ with the
same formal matrix product vector family. Its conclusion has the same shape as
the sector-comparison part of the CPSV statement: after a positive blocking,
there are sector decompositions $P$ and $Q$ with basis-of-normal-tensor data, the
blocked tensors agree with these sector decompositions at positive lengths, the
sector decompositions have the same full matrix product vector family, and the
sector multiplicities and weights match up to a permutation and nonzero phases.
The former blocked-normal-form records that encoded this boundary have been
removed from the formal development. Mathematically, the remaining hypotheses
are the comparison field for the common primitive nonzero-sector families
produced by the preceding structural theorem; the blocked-word relabeling
assertion for cyclic-sector data is derived in the formal theorem. The comparison
field is supplied with trace-preserving, primitive, and irreducible structural
data and positive bond dimensions for the produced families; it then returns
BNT block-comparison data containing the zero-tail, injectivity, normal-form
separation, and block-permutation gauge-phase comparison needed for those
produced families.
In schematic form, the conditional theorem has the boundary
\[
    \begin{aligned}
         & \operatorname{SameMPV}(A,B)
        \quad+\quad
        \mathcal C_{\mathrm{prod}}(A,B) \\
         & \hspace{4em}\Longrightarrow
        \exists p,\ \operatorname{BNTCompare}(A^{[p]},B^{[p]}),
    \end{aligned}
    \tag{Cond}
\]
where the comparison data $\mathcal C_{\mathrm{prod}}(A,B)$ contains
the produced-sector hypotheses rather than deriving them from
$\operatorname{SameMPV}(A,B)$.

These hypotheses are not written as hypotheses of Theorem~\texttt{thm1} or of
Corollary~\texttt{II\_cor2} in the cited papers. They are the formal boundary
between the canonical-form construction already proved in \Lean{} and the
compressed source argument which passes from equality of total matrix product
vectors to the BNT comparison theorem.

\section{Paper-level inputs still to be formalized}

Two parts of the named hypotheses correspond to genuine paper-level mathematics
which the sources treat as part of the theorem or as standard cited input, but
which is not yet derived in the present Lean proof without using the periodic
Fundamental Theorem, for the specific common-sector families produced by the
construction.

First, the injectivity assumptions on the produced nonzero-sector blocks
correspond to the paper's use of quantum Wielandt and block injectivity. The
sources do not assume injectivity in Theorem~\texttt{thm1}; rather, they explain
that one obtains injective or block-injective tensors after further blocking.
The formal proof must still show that this refinement holds for the final
common primitive sector families to which the overlap-rigidity comparison is
applied.
For the final produced families $P_s,Q_t$, the missing statement is not an
assumption of the paper theorem but a derivation of
\[
    \forall s,\ \exists L_s,\quad
    \operatorname{span}\{P_s(w): |w|=L_s\}=M_{D_s}(\mathbb C),
    \tag{Prod-Inj}
\]
and the corresponding statement for $Q_t$.

Second, the BNT block-comparison hypothesis contains the paper-level comparison itself,
specialized to the produced nonzero-sector families. In the source, once the
tensors are in canonical form and the bases of normal tensors have been chosen,
proportionality or equality of all matrix product vectors matches the basis
tensors up to phases and gauges. The formal theorem in \ghpr{1104} still takes
the corresponding block-comparison data as input for the specific common-sector
families: normal canonical-form data on both sides, proof that distinct sectors
are not gauge-phase equivalent, zero-tail and injectivity data, and
block-matching data. A complete proof without the periodic
Fundamental Theorem must derive all of this data from the equality of the
original tensors' matrix product vector families, after all blockings and
reindexings have been accounted for.
Let $Z_A,Z_B$ denote the zero-tail dimensions and let $\omega_A,\omega_B$
denote the weight families on the final produced sectors. For a permutation
$\sigma$, the pulled-back family is
$(\sigma^*\omega_B)_j=\omega_{B,\sigma(j)}$. The missing
comparison derivation is therefore the implication
\[
    \begin{aligned}
        \operatorname{SameMPV}(A,B)\Longrightarrow\quad
         & g_A=g_B,                                  \\
         & \exists\sigma,\phi_j,X_j,\quad
        Q_{\sigma(j)}^i=e^{i\phi_j}X_jP_j^iX_j^{-1}, \\
         & Z_A=Z_B,                                  \\
         & \omega_A=\sigma^*\omega_B.
    \end{aligned}
    \tag{Cover}
\]
for the same final produced sectors $P_j,Q_j$.

These two inputs are therefore not extra assumptions proposed by the papers. They
are parts of the paper argument that remain to be connected to the current
formal construction.

\section{Coordinate and transport conditions in the formal representation}

The other hypotheses are best understood as coordinate and transport data exposed
by the formal representation. They are mathematical assertions, but the papers do
not state them because their notation identifies the relevant objects implicitly.

The blocked-word relabeling theorem says that the nonzero part obtained by
blocking once at the common length agrees with the same sector family obtained by
first removing each period and then grouping the resulting blocks. In ordinary
paper notation both tensors are simply regarded as tensors over words of the
common blocked length. In the formal theorem, this coordinate identification is
now derived from the grouped-block cast construction.
The coordinate assertion is the equality of two descriptions of the same
blocked word:
\[
    R\!\left(A_k^{[p]}\right)(i_1,\ldots,i_p)
    =
    A_k^{i_1}\cdots A_k^{i_p}.
    \tag{Relabel}
\]

The zero-tail condition has a similar origin. The source definition of matrix
product vectors is for positive lengths, and zero blocks are invisible there.
The formal relation \leanid{SameMPV₂} also includes the empty word. A zero
block of bond dimension $D_0$ contributes nothing at positive length but
contributes $D_0$ at length zero. Hence the formal proof records a zero-tail
dimension and must either identify the two zero tails or keep the length-zero
identity separate. This is a consequence of the chosen formal convention, not a
hypothesis appearing in the paper theorem.
In formulas, for a zero block $0_{D_0}$,
\[
    \tr(0_{D_0}^{\,N})=0\quad(N>0),
    \qquad
    \tr(0_{D_0}^{\,0})=D_0.
    \tag{Zero}
\]

The transport of weights and zero-tail identities through the chosen common-sector
representatives is likewise a coordinate-transport condition. The comparison theorem must be
applied to exactly the final family of blocks after common blocking and the
canonical relabeling supplied by the structural route.
The paper moves between these descriptions without naming each transport map; the
formal proof must show that the weights, nonzero-sector tensors, and zero-tail
identities survive these changes of description.

\section{Conclusion}

The papers do not state the hypotheses of
\emph{the blocked-normal-form comparison boundary} as explicit theorem
hypotheses. The conditional theorem in \ghpr{1104} is therefore not yet the
unconditional CPSV/CPGSV statement. It records the theorem form proved by the
present derivation without the periodic Fundamental Theorem, with the remaining
paper-level inputs and formal coordinate transports made visible.

This note records a proof gap in the present Lean development rather than a
source error. The cited papers compress the canonical-form reduction, blocking,
BNT matching, and injectivity refinements in the usual paper notation. The Lean
development should keep the theorem conditional until the injectivity
refinement, the BNT block comparison for the produced sectors, and the zero-tail
and weight transports have been proved for the same final common-sector data.

\bibliographystyle{alpha}
\bibliography{references}

\end{document}
